\documentclass[11pt,a4paper]{article}

% --- Codificación e idioma ---
\usepackage[utf8]{inputenc}
\usepackage[T1]{fontenc}
\usepackage[spanish,es-tabla]{babel}
\usepackage{lmodern}

% --- Matemáticas ---
\usepackage{amsmath,amssymb,amsfonts}
\usepackage{mathtools}

% --- Layout ---
\usepackage[a4paper,margin=2.5cm]{geometry}
\usepackage{parskip}
\usepackage{microtype}

% --- Tablas ---
\usepackage{booktabs}
\usepackage{array}
\usepackage{tabularx}
\usepackage{longtable}

% --- Listas con casillas para checklist ---
\usepackage{enumitem}
\usepackage{amssymb}
\newlist{checklist}{itemize}{2}
\setlist[checklist]{label=$\square$, leftmargin=*}

% --- Listings (código MATLAB) ---
\usepackage{listings}
\usepackage{xcolor}
\definecolor{codegray}{gray}{0.95}
\definecolor{codegreen}{rgb}{0,0.5,0}
\definecolor{codeblue}{rgb}{0,0,0.7}
\lstset{
    backgroundcolor=\color{codegray},
    basicstyle=\ttfamily\footnotesize,
    breaklines=true,
    columns=fullflexible,
    commentstyle=\color{codegreen}\itshape,
    keywordstyle=\color{codeblue}\bfseries,
    frame=single,
    rulecolor=\color{black!30},
    showstringspaces=false,
    tabsize=2,
    captionpos=b,
    literate=%
        {á}{{\'a}}1 {é}{{\'e}}1 {í}{{\'i}}1 {ó}{{\'o}}1 {ú}{{\'u}}1
        {Á}{{\'A}}1 {É}{{\'E}}1 {Í}{{\'I}}1 {Ó}{{\'O}}1 {Ú}{{\'U}}1
        {ñ}{{\~n}}1 {Ñ}{{\~N}}1 {¿}{{?`}}1 {¡}{{!`}}1
        {°}{{\textdegree}}1
}
\lstdefinestyle{matlab}{language=Matlab}
\lstdefinestyle{plain}{
    language=,
    keywordstyle=\color{black},
    backgroundcolor=\color{codegray},
    basicstyle=\ttfamily\footnotesize
}

% --- Hipervínculos ---
\usepackage[hidelinks]{hyperref}

% --- Bloques destacados ---
\usepackage{mdframed}
\newmdenv[
    linecolor=black!40,
    backgroundcolor=black!3,
    linewidth=0.5pt,
    skipabove=8pt,
    skipbelow=8pt,
    innertopmargin=6pt,
    innerbottommargin=6pt
]{guia}
\newmdenv[
    linecolor=blue!40,
    backgroundcolor=blue!3,
    linewidth=0.5pt,
    skipabove=8pt,
    skipbelow=8pt,
    innertopmargin=6pt,
    innerbottommargin=6pt
]{ejemplo}

% Marcadores de calidad de ejemplo (sin emojis Unicode problemáticos)
\newcommand{\malo}{\textcolor{red!70!black}{\textbf{[MAL]}}}
\newcommand{\medio}{\textcolor{orange!80!black}{\textbf{[ACEPTABLE]}}}
\newcommand{\bueno}{\textcolor{green!50!black}{\textbf{[BIEN]}}}

% =====================================================================
\title{Guía de presentación de documentación y script\\\large{formato paper}}
\author{Proyecto: Sistema de movilidad urbana con control adaptativo topográfico\\
        PID UTN --- AMUTNPA0007734\\[4pt]
        Becado: Esteban Fidel Sian (BINID)}
\date{}

\begin{document}
\maketitle

\noindent\textbf{Objetivo de la guía:} que el avance final del becado quede formateado y documentado de modo que pueda transformarse, con mínimo trabajo adicional, en una ponencia o paper para congresos del rubro (CASE, AADECA, RPIC, ARGENCON, IEEE LATINCON, entre otros).

\hrulefill

\section*{Cómo usar esta guía}

Cada apartado incluye \textbf{preguntas guía} para auto-evaluar lo que se redacta, y \textbf{ejemplos ilustrativos} del tipo \malo{} ``así no'' / \bueno{} ``así sí''. Si el becado no puede distinguir entre los ejemplos buenos y malos, ése es el punto de la guía donde tiene que detenerse y trabajar más.

\hrulefill

\section{Filosofía general}

Un avance de beca y un paper no son lo mismo, pero deben \textbf{convivir en un único documento bien escrito}. La regla práctica es:

\begin{itemize}
    \item Estructura del paper.
    \item Profundidad y honestidad de un informe interno.
    \item Reproducibilidad de un trabajo experimental.
\end{itemize}

Esto significa que el informe del becado se redacta \textbf{como si fuera un paper corto} (6 a 10 páginas), con todas las secciones propias del formato. Lo que sobre o falte respecto a un paper publicable se ajusta luego.

\begin{guia}
\textbf{Preguntas guía generales:}
\begin{itemize}
    \item ¿Mi informe se entiende leyéndolo solo, sin que yo esté al lado explicando?
    \item ¿Si lo lee alguien que no conoce el PID, entiende el problema, lo que hice y para qué sirve?
    \item ¿Cada afirmación que hago está respaldada por una ecuación, una simulación, una figura o una cita? (Si no lo está, es opinión, y la opinión no va en un paper.)
\end{itemize}
\end{guia}

\hrulefill

\section{Estructura del documento (formato IEEE)}

\subsection{Plantilla y formato físico}

\begin{itemize}
    \item \textbf{Plantilla:} IEEE Conference (\texttt{IEEEtran} en LaTeX, o la plantilla \texttt{.docx} equivalente). Bajar de \url{https://www.ieee.org/conferences/publishing/templates.html}.
    \item \textbf{Idioma:} español (los congresos argentinos lo aceptan). Título y abstract también en inglés.
    \item \textbf{Tamaño:} 6 páginas como objetivo, 8 como máximo.
    \item \textbf{Tipografía y márgenes:} los que vienen en la plantilla, sin modificar.
    \item \textbf{Ecuaciones:} numeradas todas, alineadas a la derecha, escritas con editor matemático (no como texto plano). Variables en cursiva ($v$, $i$), vectores en negrita ($\mathbf{x}$), matrices en mayúscula negrita ($\mathbf{A}$).
    \item \textbf{Figuras:} vectoriales (\texttt{.eps} o \texttt{.pdf}) generadas desde MATLAB con \texttt{print -depsc} o \texttt{exportgraphics}. Resolución mínima 300 dpi si son raster.
    \item \textbf{Tablas:} con \texttt{booktabs} (LaTeX) o estilo limpio sin sombreados.
\end{itemize}

\subsection{Secciones obligatorias}

\begin{lstlisting}[style=plain]
Titulo
Autores y filiacion
Abstract (espanol)  +  Abstract (English)
Index Terms / Palabras clave
I.   Introduccion
II.  Marco teorico y modelo matematico
III. Representacion en variables de estado
IV.  Analisis del sistema
V.   Resultados de simulacion
VI.  Discusion
VII. Conclusiones y trabajo futuro
Agradecimientos
Referencias
\end{lstlisting}

\subsection{Contenido esperado por sección}

\subsubsection{Título}

Concreto y específico. Evitar genérico.

\begin{guia}
\textbf{Preguntas guía:}
\begin{itemize}
    \item ¿Mi título dice qué hice, sobre qué sistema, y con qué enfoque?
    \item ¿Alguien que busque ``freno regenerativo'' o ``variables de estado'' lo encontraría por palabras clave?
\end{itemize}
\end{guia}

\begin{ejemplo}
\textbf{Ejemplos ilustrativos:}
\begin{itemize}
    \item \malo{} \textit{``Estudio sobre bicicletas eléctricas''} (vago, no dice nada).
    \item \malo{} \textit{``Análisis de un sistema''} (peor todavía).
    \item \medio{} \textit{``Modelado de bicicleta eléctrica''} (correcto pero pobre).
    \item \bueno{} \textit{``Modelo en variables de estado de una bicicleta eléctrica con freno regenerativo orientado a control adaptativo topográfico''} (concreto, específico, con palabras clave).
\end{itemize}
\end{ejemplo}

\subsubsection{Abstract (\texorpdfstring{$\leq 200$}{<=200} palabras)}

En un párrafo único debe responder a: qué problema, qué se hizo, qué se obtuvo numéricamente, qué se concluyó. Sin citas. Sin acrónimos no expandidos. La versión en inglés es traducción fiel, no resumen distinto.

\begin{guia}
\textbf{Preguntas guía:}
\begin{itemize}
    \item ¿Mi abstract menciona al menos un número concreto (eficiencia, tiempo de establecimiento, error)?
    \item ¿Si alguien lee solo el abstract, sabe qué hice y qué obtuve?
    \item ¿Hay siglas sin expandir o citas? (No tiene que haber.)
\end{itemize}
\end{guia}

\begin{ejemplo}
\textbf{Ejemplos ilustrativos:}

\malo{} Abstract malo (vago, sin números, no dice qué se concluye):

\textit{``En este trabajo se estudia el modelado de bicicletas eléctricas con freno regenerativo. Se utilizan ecuaciones de variables de estado y se hacen simulaciones en MATLAB. Los resultados son satisfactorios y se discuten las posibilidades de control. Se concluye que el sistema funciona.''}

\bueno{} Abstract bueno (concreto, con números, dice qué se concluye):

\textit{``Este trabajo presenta un modelo en variables de estado de una bicicleta eléctrica con freno regenerativo, formulado con la pendiente del terreno como entrada de perturbación medible. El modelo es lineal en torno a un punto de operación de $5\,$m/s e incluye el acoplamiento electromecánico de un motor BLDC tipo hub. Se verifica que el sistema es completamente controlable y observable por criterio de Kalman. Mediante simulación en MATLAB se evalúan tres escenarios topográficos ($\theta = 0^\circ, \pm 5^\circ$) y se obtiene una eficiencia de regeneración del $47\%$ en frenados típicos urbanos, consistente con los rangos reportados en la literatura. Un análisis de sensibilidad muestra variaciones de hasta $19\%$ en el polo dominante ante cambios de $\pm 20\%$ en los parámetros del motor, lo que justifica la necesidad de un control adaptativo. El modelo es la base para el diseño del controlador adaptativo topográfico previsto en la siguiente etapa del proyecto.''}
\end{ejemplo}

\subsubsection{Index Terms}

4 a 6 palabras clave. Tomar las del PID: \textit{control adaptativo, freno regenerativo, motor brushless DC, movilidad urbana, variables de estado}.

\subsubsection{I. Introducción}

Tres bloques:

\begin{enumerate}
    \item \textbf{Contexto y motivación} (1 párrafo): movilidad eléctrica urbana en Argentina, particularidades de Paraná, recuperación de energía.
    \item \textbf{Estado del arte y trabajos relacionados} (1--2 párrafos): citar al menos 5 referencias del rubro. Usar la bibliografía disponible en el PID (Oman \& Morchin, Chau, Larminie \& Lowry, Rashid, Hendershot \& Miller).
    \item \textbf{Aporte y organización del trabajo} (1 párrafo): qué hace este paper, qué no hace, y la estructura de las secciones siguientes.
\end{enumerate}

\begin{guia}
\textbf{Preguntas guía:}
\begin{itemize}
    \item ¿Mi introducción contesta ``¿por qué importa esto?'' antes de ``¿qué hice?''?
    \item ¿Cito al menos 5 trabajos previos y los agrupo según qué hicieron, no como una lista enumerativa?
    \item ¿El último párrafo dice claramente qué hace este paper y qué no?
\end{itemize}
\end{guia}

\begin{ejemplo}
\textbf{Ejemplos ilustrativos:}

\malo{} Apertura mala (sin gancho, sin contexto, demasiado genérica):

\textit{``En este trabajo se realiza el modelado de un sistema. La movilidad eléctrica es un tema importante. Existen muchos trabajos al respecto\dots''}

\bueno{} Apertura buena (con contexto regional, motivación clara):

\textit{``La movilidad urbana eléctrica de baja escala ---bicicletas, monopatines y scooters--- ha crecido sostenidamente en la última década en Argentina, impulsada por la búsqueda de transporte personal sustentable y por la ausencia de homologación obligatoria para circular en vía pública [1]. En ciudades con topografía variable como Paraná, donde se alternan tramos planos largos con barrancas pronunciadas hacia el río, el desempeño de un vehículo eléctrico personal depende fuertemente de cuán bien el sistema de propulsión y frenado se adapte a la pendiente del momento. En particular, el frenado regenerativo es atractivo no por la magnitud de la energía recuperada ---relativamente baja en bicicleta [2]--- sino por su integración natural con un control que regule simultáneamente el frenado seguro y la recarga de la batería.''}
\end{ejemplo}

\subsubsection{II. Marco teórico y modelo matemático}

Esta es la sección donde va el desarrollo del documento Word original, \textbf{corregido} según la guía 1. Debe contener:

\begin{itemize}
    \item Diagrama de cuerpo libre de la bicicleta.
    \item Diagrama de circuito equivalente del motor BLDC.
    \item Ecuación de Newton con todas las fuerzas (aero, rodadura, gravitatoria).
    \item Ecuación eléctrica del motor (modo motor y modo generador, con convención de signos explícita).
    \item Linealización con punto de operación declarado.
\end{itemize}

\begin{guia}
\textbf{Preguntas guía:}
\begin{itemize}
    \item ¿Otra persona puede llegar a las mismas ecuaciones partiendo de mis hipótesis y diagramas?
    \item ¿Las hipótesis simplificadoras están listadas explícitamente al principio (ej.: ``se asume aire en calma, sin viento'', ``se desprecia la dinámica de la batería'')?
    \item ¿Toda variable se define la primera vez que aparece?
\end{itemize}
\end{guia}

\subsubsection{III. Representación en variables de estado}

Aquí se consolida el modelo final:

\begin{itemize}
    \item Definición de estados, entradas, salidas y perturbaciones.
    \item Matrices $A$, $B_u$, $B_d$, $C$, $D$ con dimensiones y unidades.
    \item Diagrama de bloques del sistema completo (motor + dinámica + convertidor + perturbación topográfica).
\end{itemize}

\subsubsection{IV. Análisis del sistema}

Aquí van los resultados de Kalman:

\begin{itemize}
    \item Cálculo de la matriz de controlabilidad y su rango.
    \item Cálculo de la matriz de observabilidad y su rango.
    \item Polos del sistema (analíticos, en términos de los parámetros, si se pueden simplificar; o numéricos).
    \item Discusión de qué implican estos resultados para el diseño del controlador.
\end{itemize}

\subsubsection{V. Resultados de simulación}

Mostrar las figuras desarrolladas según §2.3 de la guía 1. Cada figura \textbf{debe} referirse al texto y discutirse, no quedar suelta. Una buena estructura es: una sub-sección por escenario.

\begin{itemize}
    \item V.A. Respuesta del sistema en plano ($\theta = 0^\circ$).
    \item V.B. Comportamiento en subida ($\theta = +5^\circ$).
    \item V.C. Frenado regenerativo en bajada ($\theta = -5^\circ$). \textbf{Sección clave} del trabajo.
    \item V.D. Análisis de sensibilidad paramétrica.
    \item V.E. Balance energético.
\end{itemize}

\begin{guia}
\textbf{Preguntas guía:}
\begin{itemize}
    \item ¿Cada figura está referenciada en el texto antes de aparecer? (Ej.: ``\dots como se observa en la Fig.~3\dots'')
    \item ¿Cada subsección tiene una conclusión propia, aunque sea de una frase?
    \item ¿La sección V.C (la clave del trabajo) tiene más espacio y desarrollo que las otras?
\end{itemize}
\end{guia}

\subsubsection{VI. Discusión}

Tres preguntas a responder en orden:

\begin{enumerate}
    \item ¿Los resultados son físicamente razonables? Comparar con los rangos de Oman \& Morchin (Tabla 4.2) y con datos del fabricante del motor.
    \item ¿Qué limitaciones tiene el modelo? Mencionar: linealización local, batería ideal, convertidor en estado promediado, sin saturaciones.
    \item ¿Qué consecuencias tiene esto para el control adaptativo del PID?
\end{enumerate}

\begin{guia}
\textbf{Preguntas guía:}
\begin{itemize}
    \item ¿Soy capaz de admitir limitaciones de mi propio trabajo? (Si la sección de limitaciones es vacía o débil, el revisor desconfía del resto.)
    \item ¿La discusión hace puente con la próxima etapa o queda en el aire?
\end{itemize}
\end{guia}

\begin{ejemplo}
\textbf{Ejemplo ilustrativo --- discusión honesta vs.\ defensiva:}

\malo{} Defensiva: \textit{``El modelo predice perfectamente el comportamiento del sistema en todos los escenarios.''}

\bueno{} Honesta: \textit{``El modelo predice adecuadamente la dinámica del sistema dentro del rango $\pm 30\%$ del punto de operación elegido ($5\,$m/s). Para velocidades muy bajas ($<1\,$m/s), la linealización del término aerodinámico introduce un error superior al $25\%$, y para velocidades altas ($>8\,$m/s) la hipótesis de pequeñas señales deja de ser válida. Estos rangos se acotarán experimentalmente en la siguiente etapa del proyecto.''}
\end{ejemplo}

\subsubsection{VII. Conclusiones y trabajo futuro}

Un párrafo de conclusiones (qué se logró, en presente perfecto: \textit{``Se desarrolló\dots'', ``Se verificó\dots'', ``Se obtuvo\dots''}) y un párrafo de trabajo futuro (qué viene en la próxima etapa del PID: identificación de parámetros sobre el motor real, diseño del controlador adaptativo, validación experimental en el banco).

\subsubsection{Agradecimientos}

Mencionar a la UTN-FRP, al PID AMUTNPA0007734, al programa BINID, al director Ing.\ Katzenelson y co-director Ing.\ Maxit.

\subsubsection{Referencias}

Estilo IEEE (numéricas entre corchetes). Mínimo 10 referencias. Mezclar libros del PID, papers de IEEE Xplore y normativa cuando aplique.

\hrulefill

\section{Cómo manejar las ecuaciones}

\subsection{Editor}

En LaTeX usar entornos \texttt{equation} o \texttt{align}. En Word usar el \textbf{Editor de Ecuaciones} (no fórmulas escritas como texto plano). El borrador original tiene ecuaciones tipeadas como texto, y eso por sí solo descarta cualquier presentación seria del trabajo.

\subsection{Numeración}

Numerar todas las ecuaciones que se referencien en el texto. No numerar ecuaciones que no se referencien.

\subsection{Notación coherente en todo el documento}

Definir de una vez y mantener:

\begin{table}[h]
\centering
\begin{tabular}{l l l}
\toprule
Símbolo & Significado & Unidad \\
\midrule
$v$ & velocidad lineal de la bicicleta & m/s \\
$\omega$ & velocidad angular de la rueda & rad/s \\
$i$ & corriente del motor & A \\
$V_m$ & tensión en bornes del motor & V \\
$V_b$ & tensión nominal de la batería & V \\
$D$ & ciclo de trabajo del convertidor & --- \\
$\theta$ & ángulo de la pendiente & rad \\
$\mathbf{x}$ & vector de estados $[v;\,i]$ & --- \\
$\mathbf{A},\mathbf{B}_u,\mathbf{B}_d,\mathbf{C}$ & matrices del sistema & --- \\
\bottomrule
\end{tabular}
\end{table}

Una sola convención para todo el paper. Si se usa $\omega$ para la velocidad angular, no aparece de pronto $n$ o rpm. Si se usa $V_m$, no se usa después $V_t$ o $V(t)$.

\begin{guia}
\textbf{Preguntas guía:}
\begin{itemize}
    \item ¿Toda variable que aparece en mis ecuaciones está definida antes de usarse?
    \item ¿Uso el mismo símbolo para la misma cosa en todo el paper?
    \item ¿Las unidades acompañan a la variable la primera vez que aparece?
\end{itemize}
\end{guia}

\begin{ejemplo}
\textbf{Ejemplos ilustrativos --- ecuaciones tipeadas mal vs.\ bien:}

\malo{} Como texto plano sin editor (versión actual del borrador):
\begin{lstlisting}[style=plain]
mtot dvtdt = Fmott - Fres(t)
\end{lstlisting}

\bueno{} Con editor matemático (cómo debe verse):
\begin{equation}
m_{tot}\,\frac{dv(t)}{dt} = F_{mot}(t) - F_{res}(t)
\end{equation}
\end{ejemplo}

\subsection{Linealización}

Mostrar el punto de operación elegido y las hipótesis. Una ecuación clave es:
\begin{equation}
F_{aero}(v) = \tfrac{1}{2}\rho A C_d v^2 \approx F_{aero}(v_0) + \rho A C_d\,v_0\,(v - v_0)
\end{equation}
con la pendiente $\rho A C_d v_0$ que pasa a ser el coeficiente $B$ del modelo lineal. Eso justifica el \texttt{B} del script.

\hrulefill

\section{Cómo manejar las figuras}

\subsection{Calidad}

Generadas siempre desde MATLAB. No screenshots. Comando recomendado:

\begin{lstlisting}[style=matlab]
exportgraphics(gcf, 'figuras/fig01_polos.pdf', 'ContentType', 'vector');
exportgraphics(gcf, 'figuras/fig01_polos.png', 'Resolution', 300);
\end{lstlisting}

\subsection{Estilo uniforme}

Definir al principio del script un \texttt{set(0,\,'DefaultAxesFontSize',\,11)} y demás defaults para que todas las figuras se vean iguales. Un buen preámbulo:

\begin{lstlisting}[style=matlab]
set(0, 'DefaultAxesFontName',   'Times');
set(0, 'DefaultAxesFontSize',   11);
set(0, 'DefaultLineLineWidth',  1.5);
set(0, 'DefaultAxesGridAlpha',  0.3);
set(0, 'DefaultFigurePaperPositionMode', 'auto');
\end{lstlisting}

\subsection{Etiquetas}

\begin{itemize}
    \item Ejes siempre etiquetados con magnitud y unidad: \texttt{Velocidad v [m/s]}, no solo \texttt{v}.
    \item Leyenda dentro del gráfico cuando hay más de una curva.
    \item Título corto o sin título (preferible sin título; el caption de la figura cumple esa función en el paper).
\end{itemize}

\subsection{Caption}

El caption debe ser autocontenido: alguien que mira solo la figura tiene que entender qué representa.

\begin{guia}
\textbf{Preguntas guía:}
\begin{itemize}
    \item ¿Si tapo el cuerpo del paper y miro solo la figura y su caption, entiendo qué muestra y por qué importa?
    \item ¿El caption menciona al menos un valor numérico relevante?
\end{itemize}
\end{guia}

\begin{ejemplo}
\textbf{Ejemplos ilustrativos --- captions:}

\malo{} Mínimo (no informativo): \textit{``Fig.~3. Respuesta al escalón.''}

\medio{} Aceptable pero pobre: \textit{``Fig.~3. Respuesta del sistema al escalón en la entrada de tensión.''}

\bueno{} Bueno (autocontenido): \textit{``Fig.~3. Respuesta al escalón en la tensión del motor ($V_m = 24\,$V) con $\theta = 0^\circ$. Se observa una constante de tiempo dominante $\tau \approx 1.8\,$s asociada al modo mecánico y un transitorio rápido ($\approx 20\,$ms) asociado al modo eléctrico. La velocidad de equilibrio es de $4.8\,$m/s.''}
\end{ejemplo}

\hrulefill

\section{Cómo manejar las tablas}

Reglas mínimas:

\begin{itemize}
    \item Sin sombreados de colores.
    \item Sin bordes verticales.
    \item Bordes horizontales solo arriba, debajo del encabezado y al final.
    \item Encabezado con unidades.
    \item Caption arriba (en IEEE las tablas llevan caption arriba, las figuras abajo).
    \item Numeradas con romano (Tabla I, Tabla II), las figuras con arábigo (Fig.~1, Fig.~2).
\end{itemize}

\begin{guia}
\textbf{Preguntas guía:}
\begin{itemize}
    \item ¿Cada columna tiene unidad en el encabezado?
    \item ¿La tabla puede leerse aislada del cuerpo del texto?
\end{itemize}
\end{guia}

\begin{ejemplo}
\textbf{Ejemplos ilustrativos --- tablas:}

\malo{} Tabla mal armada:

\begin{tabular}{ll}
\toprule
Param & Valor \\
\midrule
masa & 95 \\
rad  & 0.33 \\
\bottomrule
\end{tabular}

\medskip

\bueno{} Tabla bien armada:

\textbf{TABLA I --- Parámetros del modelo nominal}

\begin{tabular}{l c c c l}
\toprule
Parámetro & Símbolo & Valor & Unidad & Fuente \\
\midrule
Masa total      & $m$    & 95   & kg & [3, p.~30] \\
Radio rueda     & $r$    & 0.33 & m  & Medido \\
Coef.\ arrastre & $C_d$  & 1.0  & ---& [3, Tabla 2.1] \\
$\dots$ & $\dots$ & $\dots$ & $\dots$ & $\dots$ \\
\bottomrule
\end{tabular}
\end{ejemplo}

\hrulefill

\section{Cómo organizar el código}

\subsection{Estructura de archivos}

\begin{lstlisting}[style=plain]
proyecto_bici_regenerativa/
|-- README.md
|-- doc/
|   |-- informe.tex (o informe.docx)
|   |-- informe.pdf
|   `-- referencias.bib
|-- src/
|   |-- modelo_bici.m            % Funcion que arma el modelo
|   |-- parametros.m             % Devuelve struct con parametros
|   |-- Modelado_V2.m            % Script principal
|   |-- Analisis_energia.m       % Balance energetico
|   `-- Sensibilidad.m           % Analisis de sensibilidad
|-- figuras/
|   |-- fig01_polos.pdf
|   |-- fig02_step_vbat.pdf
|   `-- ...
`-- data/
    `-- resultados.mat
\end{lstlisting}

\subsection{Header obligatorio en cada \texttt{.m}}

\begin{lstlisting}[style=matlab]
%==========================================================================
% Modelado_V2.m
%
% Modelo en variables de estado de una bicicleta electrica con freno
% regenerativo. Calcula matrices A, B, C, D, controlabilidad y
% observabilidad por criterio de Kalman, y simula respuestas en distintos
% escenarios topograficos.
%
% Proyecto:  PID UTN AMUTNPA0007734 - Sistema de movilidad urbana con
%            control adaptativo topografico
% Autor:     Esteban Fidel Sian
% Director:  Ing. Gustavo G. Katzenelson
% Fecha:     [completar]
% Version:   2.0
% Requisitos: MATLAB R2021a o superior, Control System Toolbox
%==========================================================================
\end{lstlisting}

\subsection{Comentarios}

\begin{itemize}
    \item Cada bloque lógico debe tener un comentario que explique \textbf{qué hace} y \textbf{por qué}, no qué \textit{operadores} ejecuta.
    \item Las unidades van en el comentario del lado de cada parámetro.
    \item Las ecuaciones implementadas se citan en el comentario respecto a una sección o ecuación del informe (\texttt{\% Ec.\ (12) del informe}).
\end{itemize}

\begin{guia}
\textbf{Preguntas guía:}
\begin{itemize}
    \item ¿Mis comentarios explican el ``por qué'' o solo el ``qué''?
    \item ¿Si borro todos los comentarios, queda código indescifrable o sigue legible?
\end{itemize}
\end{guia}

\begin{ejemplo}
\textbf{Ejemplos ilustrativos --- comentarios:}

\malo{} Comentario inútil (repite lo que ya dice el código):
\begin{lstlisting}[style=matlab]
A(1,1) = -B/m;   % asigno -B/m a la posicion (1,1) de A
\end{lstlisting}

\bueno{} Comentario útil (explica el porqué y la fuente):
\begin{lstlisting}[style=matlab]
A(1,1) = -B/m;   % termino viscoso en la dinamica de v, ec. (8) del informe
\end{lstlisting}
\end{ejemplo}

\subsection{Modularidad}

Encapsular la construcción del modelo en una función reutilizable:

\begin{lstlisting}[style=matlab]
function sys = modelo_bici(p)
% MODELO_BICI Construye el modelo en variables de estado.
%   sys = modelo_bici(p) recibe una estructura p con campos
%   m, B, r, R, L, Kt, Ke y devuelve un objeto ss.
    A = [-p.B/p.m,        p.Kt/(p.m*p.r);
         -p.Ke/(p.L*p.r), -p.R/p.L      ];
    Bu = [0; 1/p.L];
    Bd = [-9.81; 0];                          % entrada de pendiente
    C  = eye(2);
    D  = zeros(2, 2);
    sys = ss(A, [Bu Bd], C, D);
    sys.InputName  = {'V_m', 'sin(theta)'};
    sys.OutputName = {'v', 'i'};
    sys.StateName  = {'v', 'i'};
end
\end{lstlisting}

Esto permite barrer parámetros sin reescribir el código.

\subsection{Versionado}

Como mínimo guardar copias incrementales: \texttt{Modelado\_V1.m}, \texttt{Modelado\_V2.m}, etc. Lo correcto es \textbf{un repositorio git} local con commits etiquetados y un \texttt{.gitignore} que excluya archivos generados (\texttt{*.asv}, \texttt{*.mat} grandes, \texttt{figuras/*.pdf} si pesan mucho).

\subsection{Reproducibilidad}

El script principal debe ejecutarse de cabo a rabo sin intervención manual: poner los parámetros en una sola estructura, generar todas las figuras, guardar resultados. Cualquiera con MATLAB y los mismos archivos debe poder reproducir el informe.

\begin{guia}
\textbf{Pregunta guía --- el test del compañero:}

Le doy mi carpeta a un compañero de la facultad que nunca vio el proyecto. Le digo: ``ejecutá \texttt{Modelado\_V2.m}''. ¿Le sale todo? ¿Le aparecen las figuras del paper? Si me llama por chat para preguntar algo, \textbf{fallé}.
\end{guia}

\hrulefill

\section{Cómo manejar las referencias}

\subsection{Gestor}

Usar un gestor (Mendeley, Zotero, JabRef). Exportar a BibTeX si se trabaja en LaTeX, o a Word con el plugin si se trabaja en Word.

\subsection{Estilo IEEE}

Numerado entre corchetes en orden de aparición.

\begin{ejemplo}
\textbf{Ejemplos ilustrativos --- referencias:}

\malo{} Mal armada (faltan datos, formato inconsistente):

\textit{[1] Oman, libro de bicicletas eléctricas, 2006.}

\bueno{} Bien armada (formato IEEE completo):

\textit{[1] H.~Oman and W.~Morchin, \textit{Electric Bicycles: A Guide to Design and Use}. Hoboken, NJ: Wiley-IEEE Press, 2006.}

\textit{[2] K.~T.~Chau, \textit{Electric Vehicle Machines and Drives: Design, Analysis and Application}. Hoboken, NJ: Wiley-IEEE Press, 2015.}

\textit{[3] J.~Larminie and J.~Lowry, \textit{Electric Vehicle Technology Explained}, 2nd ed. Chichester: Wiley, 2012.}
\end{ejemplo}

\subsection{Mínimo de referencias y dónde citar}

\begin{itemize}
    \item Cada afirmación de hecho debe tener cita.
    \item Cada parámetro tomado de literatura debe tener cita.
    \item Modelos clásicos (Newton, Kirchhoff) no requieren cita pero sí los modelos específicos del rubro.
    \item Mínimo 10 referencias para un paper corto.
\end{itemize}

\begin{guia}
\textbf{Preguntas guía:}
\begin{itemize}
    \item ¿Estoy citando al menos los 5 libros centrales del PID (Oman \& Morchin, Chau, Larminie, Rashid, Hendershot \& Miller)?
    \item ¿Cada número de tabla, cada coeficiente, cada afirmación histórica tiene su cita?
    \item ¿Cito sin haber leído? (Si la respuesta es sí, hay que leer antes de citar.)
\end{itemize}
\end{guia}

\subsection{Citas frecuentes para este trabajo}

\begin{itemize}
    \item Oman \& Morchin para la dinámica de la bicicleta y el modelo de potencia/regeneración.
    \item Chau o Larminie \& Lowry para el contexto de vehículos eléctricos.
    \item Hendershot \& Miller o Xia para el motor BLDC.
    \item Rashid (Power Electronics Handbook) para el convertidor Buck-Boost.
    \item Yiannnis et al. (System Identification and Adaptive Control) para el marco del control adaptativo.
\end{itemize}

\hrulefill

\section{Lista de control antes de entregar}

Repasar todos estos ítems antes de dar por cerrado el informe. Si alguno falla, \textbf{no está terminado}.

\begin{checklist}
    \item El título describe específicamente lo que se hizo.
    \item Hay abstract en español y en inglés, $\leq 200$ palabras cada uno.
    \item El abstract menciona al menos un resultado numérico concreto.
    \item Las palabras clave están entre las del PID o son consistentes con el rubro.
    \item Cada ecuación está numerada y escrita con editor matemático.
    \item Cada variable se define la primera vez que aparece.
    \item La notación es consistente en todo el documento.
    \item Cada figura tiene caption autocontenido y se referencia en el texto.
    \item Cada tabla tiene caption arriba y se referencia en el texto.
    \item Hay al menos 10 referencias en formato IEEE.
    \item Cada parámetro tiene fuente.
    \item Las ecuaciones lineales están justificadas con un punto de operación explícito.
    \item Hay análisis de controlabilidad y observabilidad explícitos.
    \item Hay al menos una figura por cada escenario topográfico (plano, subida, bajada).
    \item Hay análisis energético del frenado regenerativo.
    \item Hay análisis de sensibilidad paramétrica.
    \item La discusión menciona limitaciones del modelo.
    \item Las conclusiones se cierran con trabajo futuro alineado al cronograma del PID.
    \item Los scripts MATLAB se ejecutan sin error desde cero.
    \item El repositorio tiene README con instrucciones de uso.
    \item El informe compila/se exporta a PDF sin warnings.
    \item Hay revisión ortográfica completa (en español, sin ``perdidas'' sin tilde, ``regerativo'', etc.).
\end{checklist}

\begin{guia}
\textbf{Preguntas guía finales antes de mandar:}
\begin{itemize}
    \item ¿Lo leí en voz alta de principio a fin? (Es la mejor forma de detectar frases mal escritas.)
    \item ¿Lo dejé descansar 24 hs y lo releí con ojos frescos? (Casi siempre se encuentran errores nuevos.)
    \item ¿Lo leyó al menos una persona ajena al proyecto?
\end{itemize}
\end{guia}

\hrulefill

\section{Camino hacia el congreso}

Una vez aprobado el informe interno, transformarlo en envío de congreso requiere:

\begin{enumerate}
    \item \textbf{Identificar el congreso objetivo.} Sugerencias para 2026/2027 en Argentina: AADECA (Asociación Argentina de Control Automático), RPIC (Reunión de Procesamiento de la Información y Control), ARGENCON, CASE (Congreso Argentino de Sistemas Embebidos).
    \item \textbf{Adaptar la plantilla} a la del congreso (suelen ser variantes de IEEE).
    \item \textbf{Reducir o ampliar} según el límite de páginas del CFP (Call For Papers).
    \item \textbf{Pedir revisión interna} al director y co-director del PID.
    \item \textbf{Submission.} Subir antes del deadline. Conservar el comprobante.
    \item \textbf{Defensa.} Si es aceptado, preparar presentación de 15 minutos: una diapositiva por sección del paper, foco en la sección V (resultados).
\end{enumerate}

\begin{guia}
\textbf{Preguntas guía para la defensa:}
\begin{itemize}
    \item ¿Puedo explicar mi trabajo en 1 minuto a alguien que no conoce el tema? (El elevator pitch.)
    \item ¿Puedo explicarlo en 5 minutos a alguien del rubro? (La introducción de la presentación.)
    \item ¿Puedo defender cada decisión técnica que tomé? (La sesión de preguntas.)
\end{itemize}
\end{guia}

\hrulefill

\section{Consejo final}

Un buen avance de beca no es un paper publicable, pero \textbf{un buen avance se nota inmediatamente porque no requiere reescribirse para convertirse en uno}. Si al final del trabajo el director pide ``vamos a publicarlo'' y la respuesta es ``necesito 2 meses para reescribirlo'', el avance no estuvo bien hecho. Si la respuesta es ``lo adaptamos en 2 semanas a la plantilla del congreso'', el avance estuvo bien hecho.

El objetivo de esta guía es que el segundo escenario sea el resultado real.

\begin{guia}
\textbf{Pregunta guía de cierre:}

Si dentro de 5 años retomo este trabajo (o lo retoma otro becado), ¿el material que dejé hoy es suficiente para que la continuación arranque sin reconstruir lo hecho? Esa es la prueba real de un avance bien terminado.
\end{guia}

\end{document}
